\section{A five-cycle shore}
\label{sec:five-cycle-shore}

The protected-set form of the small-cut lemma makes the added-edge inequalities unconditional.
This gives a fixed deficit for the smallest cyclic five-cut shore.

\begin{lemma}[Unconditional augmented-path inequalities]
\label{lem:five-pole-unconditional-augmentation}
Fix $\varepsilon>0$.
Let $Y$ be a normalized five-pole of sufficiently large rank $k$.
Suppose a cubic completion of $Y$ has no cyclic edge cut of size at most four.
Normalize its terminal counts by $2^k$.
Then every pair satisfies $a_e\le1/2+\varepsilon$.
For every four-port set $S$ and each pairing $\alpha=ij|lm$ of $S$,
\[
 a_{ij}+U_S-u_{S,\alpha}\le1+2\varepsilon,
 \qquad
 a_{lm}+U_S-u_{S,\alpha}\le1+2\varepsilon.
\]
\end{lemma}

\begin{proof}
The ordinary path bound follows from the fixed-density small-cut theorem.
Add the edge $lm$ to $Y$.
This gives a normalized three-pole of rank $k+1$.
Its $i$--$j$ path count equals
\[
 2^k\bigl(a_{ij}+U_S-u_{S,\alpha}\bigr).
\]
If the first augmented inequality failed, this new pole would have terminal-path density greater than $1/2+\varepsilon$.
Apply the protected-set small-cut theorem with both endpoints of the added edge protected.
It supplies two vertex-disjoint cyclic sets, each with boundary at most four, that avoid those endpoints.
Both cycles and their boundaries are therefore unchanged when the added edge is removed and the original completion is restored.
This contradicts the hypothesis on that completion.
Exchange the two pairs for the other inequality.
All augmentation patterns form a fixed finite list, so one rank threshold suffices.
\end{proof}

\begin{theorem}[A five-cycle cut]
\label{thm:five-cycle-shore-sharp-order}
There is an absolute constant $K$ such that the following statement holds.
Let $G$ be a cubic graph obtained by joining a five-cycle to a normalized five-pole $Y$ of rank $k$.
Put $r=k+5$.
Then
\[
 c(G)\le2^{r-1}+Kr2^{r/2}.
\]
If $k$ is sufficiently large and $G$ has no cyclic cut of size at most four, then $c(G)<2^{r-1}$.
\end{theorem}

\begin{proof}
If $G$ has a cyclic cut of size at most four, apply Theorem~\ref{thm:sharp-order-four-cut-error}.
Otherwise assume $k$ is sufficiently large for Lemma~\ref{lem:five-pole-unconditional-augmentation}.
Label the five ports in their cyclic order on the five-cycle.
Write $S_{\rm edge}$ for the sum of the five path densities joining consecutive ports.
Write $S_{\rm diag}$ for the sum of the other five path densities, and put $S=S_{\rm edge}+S_{\rm diag}$.

Every terminal pair in the five-cycle has two paths.
In each four-port sector it has two pairing forests, one in each noncrossing pairing.
The alternating pairing is absent.
Let $A$ be the sum of the five alternating-pairing masses of $Y$, and let $R=\sum_{|T|=4}U_T$.
The crossing contribution, divided by $2^{k+1}$, is exactly
\[
 L=S+\frac12(R+A).
\]

For a noncrossing pairing $\alpha=ij|lm$ in a sector $T$, put $M_\alpha=\max(a_{ij},a_{lm})$.
The two augmented inequalities imply
\[
 M_\alpha+U_T-u_{T,\alpha}\le1+2\varepsilon.
\]
Sum this over the two noncrossing pairings of $T$.
Their two masses sum to $U_T-u_{T,\rm alt}$, so
\[
 M_{\alpha_1}+M_{\alpha_2}+U_T+u_{T,\rm alt}\le2+4\varepsilon.
\]
Across all five sectors, each consecutive pair occurs in three noncrossing pairings.
Each diagonal pair occurs in one.
Using $\max(x,y)\ge(x+y)/2$ gives
\[
 R+A\le10+20\varepsilon-\frac12(3S_{\rm edge}+S_{\rm diag}).
\]
Consequently,
\[
 L\le5+10\varepsilon+\frac14S_{\rm edge}+\frac34S_{\rm diag}
 \le\frac{15}{2}+15\varepsilon.
\]

Take $\varepsilon=1/100$.
Then $L\le7.65$.
The asymptotic cycle estimate gives $c(Y)\le(9/16)2^k$ for sufficiently large $k$.
The five-cycle contributes one internal cycle.
Hence
\[
 \frac{c(G)}{2^{k+1}}
 \le7.65+\frac9{32}+2^{-k-1}<8
\]
for sufficiently large $k$.
This proves the strict half bound in this case.
The remaining bounded ranks are absorbed by increasing the absolute constant $K$.
\end{proof}

The proof uses the exact five-cycle pairing vector $(1/2,1/2,0)$ after normalization.
The alternating pairing receives the larger gluing coefficient because it is absent on the five-cycle shore.
